\chapter{Supplementary Sweep Results}

This appendix provides additional hyperparameter sweep visualisations, sweep budgets, and the documented compute environment for runs whose quantitative findings appear in the main report's \textbf{Chapter~7 (Evaluation/Results)}. It does not reproduce those headline results, which are argued there; the Implementation chapter (Chapter~6) describes code structure only. Operational commands appear in the User Guide and Source Code appendices.

\section{Sweep Configuration Space}

The sweep budgets are dataset-specific. This was an explicit practical decision:
the MIMIC run was computationally expensive in wall-clock time, so the
Challenge 2012 profile used a smaller full-grid budget to keep the overall
experimental campaign feasible.

\begin{table}[htbp]
\centering
\small
\begin{tabular}{p{0.23\textwidth}p{0.36\textwidth}p{0.36\textwidth}}
\toprule
\textbf{Setting} & \textbf{MIMIC full profile} & \textbf{Challenge 2012 full profile} \\
\midrule
$n_t$ values & $\{1, 2, 3, 5, 8, 13, 21, 34\}$ ($8$ values) & $\{1, 2, 3, 5, 8, 13, 21\}$ ($7$ values) \\
$n_{\text{noise}}$ values & $\{1, 2, 3, 5, 8, 13, 21, 34\}$ ($8$ values) & $\{1, 2, 3, 5, 8, 13, 21\}$ ($7$ values) \\
Grid size & $8 \times 8 = 64$ cells & $7 \times 7 = 49$ cells \\
Max training rows & $100{,}000$ & $100{,}000$ \\
Synthetic samples per cell & $100{,}000$ & $50{,}000$ \\
Privacy metric row cap & $2{,}000$ & $2{,}000$ \\
\bottomrule
\end{tabular}
\caption{Dataset-specific full-profile sweep budgets from \code{src/lexisflow/config/datasets.py}.}
\label{tab:dataset-sweep-budgets}
\end{table}

Unless otherwise stated, sweeps in this appendix use the HS3F backbone and three
trajectory sampling seeds per cell
(\code{TSTR\_TRAJECTORY\_SAMPLING\_SEEDS = (42, 11, 50)}).

\section{XGBoost Configuration}

The shared sweep XGBoost parameter block (\code{SWEEP\_XGB\_PARAMS}) is:

\begin{table}[htbp]
\centering
\begin{tabular}{ll}
\toprule
\textbf{Parameter} & \textbf{Value} \\
\midrule
\code{n\_estimators} & 100 \\
\code{max\_depth} & 6 \\
\code{learning\_rate} & 0.1 \\
\code{subsample} & 0.8 \\
\code{colsample\_bytree} & 0.8 \\
\bottomrule
\end{tabular}
\caption{Shared XGBoost hyperparameters used in sweep training.}
\label{tab:xgb-params}
\end{table}

These parameters were not tuned as part of the sweep; they represent a single fixed operating point. Parallelism is controlled separately by the sweep CLI \code{--n-jobs} (default \code{8}), and model random state is fixed in the sweep factory for reproducibility.

\section{Compute Environment for Reproducibility}

All reported sweeps were executed on a local Apple Silicon workstation. The
machine characteristics are listed for transparency because runtime was a
material constraint in sweep design:

\begin{table}[htbp]
\centering
\begin{tabular}{ll}
\toprule
\textbf{Component} & \textbf{Specification} \\
\midrule
CPU & Apple M1 Max \\
Logical CPU cores & 10 \\
Memory & 32 GB RAM \\
Operating system & macOS 26.4 (Darwin 25.4.0, arm64) \\
Typical sweep parallelism & \code{--n-jobs 8} (varied by run) \\
\bottomrule
\end{tabular}
\caption{Hardware and software environment used for the reported sweeps.}
\label{tab:sweep-compute-env}
\end{table}

The 32~GB figure records the machine used for Chapter~7; it does not by itself establish that the same end-to-end sweep would complete on a smaller-RAM workstation. That is separate from NFR4, which the main report verifies as a \emph{software} requirement (iterator/streaming/caching paths used on the full sweep), not as a proof of minimum physical RAM.
